\documentclass[12pt]{article}
\usepackage[margin=1in]{geometry}
\usepackage{setspace}
\usepackage{fancyhdr}
\usepackage{hyperref}

% Set double spacing
\setstretch{1}

% Header with last name and page number
\pagestyle{fancy}
\fancyhf{}
\rhead{Vaught \thepage}

% Remove the header line
\renewcommand{\headrulewidth}{0pt}

\begin{document}

% Upper left corner information
\noindent
J.C. Vaught \\
Professor Emily Allen \\
MUSC 113 \\
4 Jan. 2025 \\

\vspace{0.5in}

% Centered title
\begin{center}
\textbf{Journal \#1, Chapter 1 and 2 Playlists}
\end{center}

\vspace{0.5in}

% Begin main content
From the selection of songs I could have chosen for Chapter 1, I decided to listen to Judy Garland’s recording of “Over the Rainbow” (the version arranged by Victor Young and His Orchestra) and played it a few times, determined to find any nods/homages being paid to the 'blues' genre. Of course, I first heard the fullness of the orchestral arrangement, with strings and wind instruments sound swirling around Judy’s voice. Not exactly what I had come to expect from a blues, something I associate with a primarily guitar and harmonica setup.

I then knew I had to formally define what a blues song was before I could answer the prompt for the Journal.

\subsection*{Attempting to Define the Blues}

When I attempted to define the Blues using the textbook for this class, I found it too limiting. Seeking a deeper understanding, I turned to Michael V. Uschan’s \textit{The History of the Blues}, a book in a series by Cengage that categorizes the history of different music genres. This book offered a more \textit{\textbf{knowledgeable}} perspective on the Blues, making my reading experience more useful.

Uschan provides a foundational perspective on the Blues, stating, “The blues was the first distinctly American music because it included rhythms and musical elements derived from African music, such as call and response and blue notes, that were foreign to Western music” (Uschan 10). This highlights how the Blues not only emerged uniquely but also laid the groundwork for numerous modern music genres.

In the early phases of the Blues, Uschan describes the primitive instruments: “The most primitive [banjos] were made of hollowed-out gourds strung lengthwise with thin strips of dried animal skin attached to a wooden neck” (Uschan 17). The use of everyday objects like jugs and washboards as instruments shows how the Blues was originally created and its lasting influence on the genre.

As the Blues evolved into the `Classic Blues Era' of the 1920s, so did its instrumentation. Uschan notes, “Smith was accompanied by five musicians who played piano, trumpet, trombone, clarinet, and violin” (Uschan 28). This shift from solo performances to ensemble settings allowed for more complex and richer musical arrangements, marking a significant development in the Blues.

The transition to `Chicago Blues' in the late 1930s and early 1940s illustrates how urbanization influenced the genre: “Bluesmen began adding other instruments to their solo guitar playing to provide a richer musical background to their vocals, including drums, pianos, harmonicas, and wind instruments” (Uschan 49). This adaptation was crucial for thriving in city nightclubs, where amplified sounds became necessary.

The Blues is deeply emotional, a reflection of its origins. Uschan explains, “The deep, pervasive current of melancholy that characterizes many blues songs arose from the unfair, often brutal treatment African Americans experienced in the United States” (Uschan 2). The universal themes of love, loss, and hardship make the Blues timeless, allowing listeners to connect with the music on a personal level.

Structurally, the Blues is distinctive with its twelve-bar format and AAB lyrical structure. Uschan explains, “The lyrics will be in AAB form, which means that the first line of the song will be repeated, with the third line acting as a sort of answer to the previous repeated lines” (Uschan 34). This repetitive and conversational style facilitates improvisation and reinforces the storytelling aspect of the Blues, making each song a personal narrative.

\subsection*{Analysis of ``Over the Rainbow"}

After refining my understanding of the blues, I revisited Garland’s rendition of “Over the Rainbow” with a sharper lens, seeking to identify any homages it might pay to the genre. One striking feature was Garland’s vocal line, where a slight “warbling” on sustained notes stood out. This gentle vibrato, dipping and soaring delicately, bore a resemblance to the blues’ tradition of bending notes for emotional expression. However, this element, while evocative, felt more aligned with the polished aesthetics of Broadway vocalists than with the rawer, grittier tones typical of blues singers in my opinion.

To delve deeper, I analyzed the spectrogram of the recording, observing subtle sinusoidal variations in pitch on these sustained notes. While reminiscent of the purposeful pitch bending in blues, it’s also a staple of theatrical/pop singing. Garland’s warble, though undeniably lovely, leaned more towards the refined rather than the visceral edge that defines the blues.

Structurally, the song diverges from blues conventions. Instead of the repetitive 12-bar progression central to blues, “Over the Rainbow” adheres to a classic AABA format, with two verses (A), a contrasting bridge (B), and a return to another verse (A). This framework, emblematic of early 20th-century pop ballads, contrasts sharply with the cyclical chord patterns and call-and-response phrasing integral to blues.

Culturally and lyrically, the song ventures even further from the blues tradition. While blues often captures personal hardship and raw emotion, “Over the Rainbow” is a tale of hopeful escapism, conjuring dreamlands and skies of blue. Its cinematic optimism lacks the intimate, lived storytelling that grounds the blues in authentic experience.

Yet, despite its divergence from the blues blueprint, the song pays subtle homage to the genre through its emotional sincerity and Garland’s nuanced vocal delivery. The warbling vibrato, while more polished, echoes the blues’ spirit of heartfelt expression. These faint nods, though fleeting, reveal the interconnectedness of musical traditions, offering a fascinating glimpse into how even a Tin Pan Alley classic like “Over the Rainbow” can brush against the edges of the blues world.

\subsection*{Analysis of ``Blueberry Hill"}

When determining which song to write about in Chapter 2, I had a much harder time since all the songs sounded especially good, particularly when compared to Chapter 1 where I was not particularly a fan of the music.

I decided to use ChatGPT to pick one for me with the most `homage' to the blues genre, and it selected ``That's All Right'' or ``Blueberry Hill.'' The ``Blueberry Hill'' song was much more interesting to me since it has a very prominent piano at the beginning of the song and it remains a prominent instrument, and the piano/keyboards are my favorite instrument so I decided to analyze the song ``Blueberry Hill'' by Fats Domino. Elvis’s song was not nearly as captivating; something about the basic guitar and bass just doesn't draw me in.

I think this song is a great example of how popular music has taken elements from blues, providing an avenue for the genre to reach more audiences. The song is reminiscent of blues with its emotion-driven vocal delivery, slower tempo, and a 12-bar blues structure that forms the backbone of the melody.

One of my favorite parts of the song is the piano introduction. The piano's sound in the opening has a unique sound profile that doesn’t immediately lend itself to a piano, giving it a distinct character that captures my attention. As the song goes on, the piano becomes more recognizable in the bluesy role, providing rhythm and depth. The skillful playing keeps the instrument at the heart of the track, a common feature of blues where the piano often serves as both a rhythmic and melodic anchor.

The instrumentation also nods to blues roots while expanding its sound palette. Around 1:20, a brass instrument, possibly a trumpet, adds a layer of richness and gives the song a swing-like feel. This subtle addition reflects the blues tradition of incorporating horns to enhance emotional expression. The drum set, heard at the start, adds a steady, syncopated beat with cymbals punctuating the rhythm, a common technique in blues and early rock 'n' roll to drive the song forward while maintaining a relaxed, laid-back feel.

Throughout the song, Domino’s soulful vocals take center stage to me, emphasizing the lyrics with heartfelt expression. This focus on vocal prominence is a key aspect of blues music, where the singer conveys the song’s emotion and narrative directly to the listener.

\newpage

\begin{thebibliography}{9}
\bibitem{uschan}
Uschan, Michael V. \textit{The History of the Blues}. Cengage, 2013.
\bibitem{garland} Garland, Judy. ``Over the Rainbow.'' Arranged by Victor Young and His Orchestra, 1939.
\bibitem{domino} Domino, Fats. ``Blueberry Hill.'' Specialty Records, 1956.
\end{thebibliography}

\end{document}
